\subsection{结构的推导}

设 \(\Sigma\) 是理论 \(\mathfrak{G}\) 中的一种结构，它有 \(n\) 个主基集 \(x_{1},\ldots ,x_{n}\) 和 \(m\) 个辅助基集 \(\mathrm{A}_1,\dots ,\mathrm{A}_m\) . 设 \(s\) 是种类 \(\Sigma\) 的一般结构， \(\mathbf{T}\) 是一个具有 \(n + m\) 个项的阶梯构造方案。若项 \(\mathrm{V}\{x_1,\dots,x_n,s\}\) 除了所在理论的常量外不含其他字母（这里的理论为 \(\mathfrak{G}_{\Sigma}\) ），则在它满足下述条件时，称它对 \(s\) 是内在的，其类型为 \(\mathrm{T}(x_1,\dots,x_n,\mathrm{A}_1,\dots,\mathrm{A}_m)\) ：

\begin{enumerate}
\def\labelenumi{(\arabic{enumi})}
\item
  该关系 \(\mathbf{V}\{x_1, \ldots, x_n, s\} \in \mathbf{T}(x_1, \ldots, x_n, A_1, \ldots, A_m)\) 是定理于 \(\mathfrak{G}_{\Sigma}\) ;
\item
  让 \(\mathfrak{G}_{\Sigma}^{\prime}\) 为由以下公理添加到 \(\mathfrak{G}_{\Sigma}\) 的公理''所得到的理论 \(f_{i}\) 是 \(x_{i}\) 到上 \(y_{i}''\) ( \(1 \leqslant i \leqslant n\) ) 的双射（其中字母 \(y_{i}\) 并且 \(f_{i}\) 彼此不同，且不同于 \(\mathfrak{G}_{\Sigma}\) ，用于 \(1 \leqslant i \leqslant n\) 的常数；如果 \(s'\) 是通过 \(s\) 进行传输所得到的结构 \((f_{1}, \ldots, f_{n})\) （第5号），然后
\end{enumerate}

\[
\mathrm{V} \left\{y _ {1}, \dots , y _ {n}, s ^ {\prime} \right\} = \langle f _ {1}, \dots , f _ {n}, \operatorname{Id} _ {1}, \dots , \operatorname{Id} _ {m} \rangle^ {\mathrm{T}} (\mathrm{V} \left\{x _ {1}, \dots , x _ {n}, s \right\})
\]

是定理于 \(\mathfrak{C}_{\Sigma}^{\prime}\) .

在结构种类的理论中，人们被引导去定义的大多数术语都是内在术语。

让 \(\Theta\) 是理论中另一种结构种类 \(\mathfrak{G}\) ，在 \(r\) 主基集 \(u_{1},\ldots,u_{r}\) ，以 \(p\) 辅助基集 \(\mathrm{B}_{1},\ldots,\mathrm{B}_{p}\) ，并且让 \(t\in\mathrm{T}(u_{1},\ldots,u_{r},\mathrm{B}_{1},\ldots,\mathrm{B}_{p})\) 是……的典型刻画 \(\Theta\) （第4号）。那么，一个 \(r+1\) 项 \(\mathrm{P},\mathrm{U}_{1},\ldots,\mathrm{U}_{r}\) 的系统，对于 \(s\) ，并且使得 \(\mathrm{P}\) 是属于物种 \(\Theta\) 的、在 \(\mathrm{U}_{1},\ldots,\mathrm{U}_{r}\) 上的一个结构（在理论 \(\mathfrak{G}_{\Sigma}\) 中），被称为物种 \(\Theta\) 的结构由物种 \(\Sigma\) 的结构演绎而得的一个演绎程序。滥用语言时，术语 \(\mathrm{P}\) 单独常被称为演绎程序。

让 \(\mathfrak{G}'\) 为比……更强的理论 \(\mathfrak{G}\) 。如果 \(\mathcal{G}\) 是 \(\mathfrak{G}'\) 中物种 \(\Sigma\) 在 \(\mathrm{E}_1, \ldots, \mathrm{E}_n\) 上的一个结构，那么 \(\mathrm{P}\{\mathrm{E}_1, \ldots, \mathrm{E}_n, \mathcal{S}\}\) 是物种 \(\Theta\) 在这 \(r\) 个集合 \(\mathrm{F}_j = \mathrm{U}_j\{\mathrm{E}_1, \ldots, \mathrm{E}_n, \mathcal{S}\}(1 \leqslant j \leqslant r)\) 上的一个结构，称为由 \(\mathcal{S}\) 通过过程P，或从属于 \(\mathcal{G}\) 。假设项P， \(\mathrm{U}_1, \ldots, \mathrm{U}_r\) 对……是内在的 \(s\) 此外还蕴含以下判据：

CST6。 让 \((g_{1}, \ldots, g_{n})\) 是……的一个同构 \(E_{1}, \ldots, E_{n}\) ，配备有物种 \(\Sigma\) ，映到 \(E_{1}^{\prime}, \ldots, E_{n}^{\prime}\) 的结构，配备有结构 \(S^{\prime}\) 同一物种的。如果 \(U_{j}\) 是类型 \(\mathfrak{P}(T_{j})\) 的严格递增映射，令

\[
h _ {j} = \left\langle g _ {1}, \dots , g _ {n}, \operatorname{Id} _ {1}, \dots , \operatorname{Id} _ {m} \right\rangle^ {\mathrm{T}_j} \quad (1 \leqslant j \leqslant r),
\]

并令 \(\mathbf{F}_j' = \mathbf{U}_j\{\mathbf{E}_1', \ldots, \mathbf{E}_n', \mathcal{S}'\} (1 \leqslant j \leqslant r)\) 。然后 \((h_1, \ldots, h_r)\) 是……的一个同构 \(\mathbf{F}_1, \ldots, \mathbf{F}_r\) 到上 \(\mathbf{F}_1', \ldots, \mathbf{F}_r'\) 当这些集合系统被赋予物种结构时 \(\Theta\) 由...推导出 \(\mathcal{S}\) 并且 \(\mathcal{S}'\) 分别通过过程P。

显然，项 \(x_{1}, \ldots, x_{n}\) 对……是内在的 \(s\) 在许多情况下，项 \(U_{1}, \ldots, U_{r}\) 是某些字母 \(x_{1}, \ldots, x_{n}\) ；物种的结构 \(\Theta\) 由...推导出 \(s\) 通过该程序 \(P\) 随后被称为底层结构 \(s\) .

\minorheading{示例}

* (1) 拓扑群结构的物种具有一个主基集A，没有辅助基集，相应的典型结构是一个对 \((s_1, s_2)\) ( \(s_1\) 是A上合成律的图，而 \(s_2\) 是A拓扑中的开集族；参见《一般拓扑学》第三章第1节）。每一项 \(s_1\) , \(s_2\) 都是一个推导过程，分别提供拓扑群结构所依据的群结构和拓扑。 \((s_1, s_2)\) .

同样地，从向量空间结构可以推导出底层的交换群结构。从环结构可以推导出底层的交换群结构和（乘法）半群结构。从可微流形结构可以推导出底层的拓扑，等等。

\begin{enumerate}
\def\labelenumi{(\arabic{enumi})}
\setcounter{enumi}{1}
\tightlist
\item
  复数域C（相应地，实数域R）上向量空间结构的物种具有一个主基集E，一个等于C（相应地，R）的辅助基集，以及典型刻画
\end{enumerate}

\[
s _ {1} \in \mathfrak {P} ((\mathbf {E} \times \mathbf {E}) \times \mathbf {E}) \quad \text { 且 } \quad s _ {2} \in \mathfrak {P} ((\mathbf {C} \times \mathbf {E}) \times \mathbf {E})
\]

\[
(\text { 分别 } \quad s _ {1} \in \mathfrak {P} ((E \times E) \times E) \quad \text { 且 } \quad s _ {2} \in \mathfrak {P} ((R \times E) \times E)).
\]

对 \((s_{1}, s_{2} \cap ((\mathbf{R} \times \mathbf{E}) \times \mathbf{E}))\) 是从复数域C上的向量空间结构推导实数域R上的向量空间结构的过程（“将标量域限制为R”）。*

\begin{enumerate}
\def\labelenumi{(\arabic{enumi})}
\setcounter{enumi}{2}
\tightlist
\item
  假设 \(\Theta\) 具有与 \(\Sigma\) 相同的主基集和辅助基集，以及相同的典型刻画。此外，如果 \(\Sigma\) 的公理蕴含（在 \(\mathfrak{C}\) ) 公理 \(\Theta\) ，显然该项 \(s\) 是从物种结构 \(\Theta\) 到物种结构的演绎程序 \(\Sigma\) . \(\Theta\) 则被称为比……更贫乏 \(\Sigma\) ，以及 \(\Sigma\) 比……更丰富 \(\Theta\) 。每个属于物种 \(\Sigma\) 的结构，在理论 \(\mathfrak{C}'\) 该理论强于 \(\mathfrak{C}\) ，则也是物种 \(\Theta\) 的一个结构。例如，全序集结构的物种（通过取序结构公理（第4节，例1）与关系
\end{enumerate}

的合取作为公理而得到）比序结构的物种更丰富。*交换群结构的物种比群结构的物种更丰富。紧致拓扑空间结构的物种比拓扑结构的物种更丰富，等等。* \(s \cup \overline{s} = A \times A\) *（4）当

各自是群结构（相应地，环结构）的物种时，在代数中定义了一种演绎程序，它将每个群结构（相应地，环结构）与其中心上的群结构（相应地，环结构）相关联。当 \(\Sigma\) 并且 \(\Theta\) 是域K上向量空间结构的物种，且 \(\Sigma\) 是K上代数结构的物种时，定义了演绎程序，将K上的每个向量空间与其张量代数或外代数相关联。在本系列的后续部分，我们还会遇到许多其他例子。* \(\Theta\) 7. 等价的结构物种

注记。当P是一个“q元组”时 \((P_{1}, \ldots, P_{q})\) ，也称这些项 \(P_{1}, \ldots, P_{q}\) 构成一个物种结构的演绎过程 \(\Theta\) 到物种结构的演绎程序 \(\Sigma\) .

